\chapter{Study Methodology}\doublespacing
\label{Chapter3}
\lhead{Chapter III. \emph{Study Methodology}}

\section{Introduction}
This chapter describes the methodology used to implement Magnavis and to compare detection schemes on labelled observatory data. The emphasis is on a reproducible offline protocol that holds preprocessing and evaluation grids constant across models.

\section{Methodology}
Figure~\ref{fig:figure3_1} summarises the pipeline. Time series are read from CSV exports or from the project database. The headless application loads a pretrained recurrent model, produces one-step-ahead forecasts, forms absolute residuals, and applies the anomaly rule. Point-level counts (true positive, false positive, true negative, false negative) are computed on seconds where manual ground truth and predictions are both defined.

\begin{figure}[!ht]
    \centering
    \begin{adjustbox}{max width=0.85\textwidth, center}
        \begin{tikzpicture}[node distance=1.4cm]
            \node (csv) [startstop] {Magnetic CSV / database};
            \node (app) [process, below of=csv] {Magnavis \texttt{app.py}};
            \node (pred) [process, below of=app] {GRU / LSTM predictor};
            \node (res) [process, below of=pred] {Residual $|B_t - \hat{B}_t|$};
            \node (det) [process, below of=res] {EWMA threshold ($\mu + k\sigma$)};
            \node (met) [process, below of=det] {Metrics vs.\ manual ground truth};
            \draw [arrow] (csv) -- (app);
            \draw [arrow] (app) -- (pred);
            \draw [arrow] (pred) -- (res);
            \draw [arrow] (res) -- (det);
            \draw [arrow] (det) -- (met);
        \end{tikzpicture}
    \end{adjustbox}
    \caption{End-to-end Magnavis pipeline for offline benchmark evaluation.}
    \label{fig:figure3_1}
\end{figure}


\section{Benchmark Protocol}
The primary evaluation bundle uses a \textbf{zero-historic, predict-only} setting: no leading training minutes inside the session, pretrained weights used as shipped, and historic context set to zero unless stated otherwise. Sensors OBS2\_1--OBS2\_3 are evaluated on a common per-second grid. The threshold multiplier $k$ is swept over $\{1,2,3,4,5\}$.

Three ground-truth layouts are used:
\begin{itemize}
    \item \textbf{Feb-13 2026:} short shot-duration windows ($\approx 4.4\%$ positive seconds).
    \item \textbf{Apr-27 2026:} long sustained-offset windows ($\approx 94.5\%$ positive seconds).
    \item \textbf{Synthetic:} flat magnetics with the long-window label geometry ($\approx 50\%$ positive).
\end{itemize}

Additional runs compare open-loop and closed-loop autoregressive inference on the Apr-27 file by rolling the input window on predicted values (\texttt{PREDICTOR\_AR\_CLOSED\_LOOP=1}).

\section{Models and Algorithms}
\textbf{Statistical baselines} forecast the series by EWMA, median filtering, or Savitzky--Golay smoothing; residuals are formed against these forecasts as for the neural models.

\textbf{Pretrained GRU and LSTM} use sliding windows of length $W=15$ with cyclic time-of-day features. Checkpoints were trained offline on a December 2025 OBS2 archive.

\textbf{Attention bidirectional LSTM} is trained only on the first 62 minutes of each evaluation CSV and then used to forecast the remainder, representing on-site adaptation.

The detector flags a second when $e_t > \mu_t + k \sigma_t$, where $\mu_t$ and $\sigma_t$ come from an EWMA of past $|e|$.

\section{Summary}
The methodology combines a fixed ingestion path, recurrent forecasting, residual-based detection, and a labelled benchmark grid. Chapter~\ref{Chapter4} details data handling; Chapter~\ref{Chapter5} reports results.
